\documentclass{article}
\usepackage[utf8]{inputenc}
\usepackage{hyperref}
\usepackage{xurl} % For better URL breaking
\usepackage{geometry}
\usepackage{fancyvrb} % For verbatim environments
\usepackage{enumitem} % For customizing list environments

\geometry{a4paper, margin=1in}
\title{LLM-Security-Lab-2026}
\author{} % Author can be added if desired
\date{} % Removes the date

\begin{document}

\maketitle
\section*{🛡️ LLM Security Lab: Prompt Injection Defenses}

In this lab, you will explore the fascinating and critical area of Large Language Model (LLM) security, specifically focusing on \textbf{jailbreaking} techniques. Your task is to build and implement robust defenses to protect a chatbot from malicious user inputs.

\subsection*{🎯 Your Goal}

Your primary goal is for this lab is to play with jailbreaks and, more importantly, to develop effective countermeasures. You will be tasked with modifying the \texttt{defense.py} file to sanitize user inputs and prevent the chatbot from revealing sensitive information or behaving in unintended ways.

\subsection*{🚀 Getting Started}

Follow these steps to set up your local development environment and get the chatbot running:

\subsubsection*{Prerequisites}

Before you begin, ensure you have the following installed:

\begin{enumerate}[label=\arabic*.]
    \item \textbf{Docker Desktop}: This lab uses \texttt{docker-compose} to manage the application and its dependencies.
    \begin{itemize}
        \item Download and install Docker Desktop from \url{https://www.docker.com/products/docker-desktop/}.
    \end{itemize}
    \item \textbf{Ollama}: We will be running a local LLM using Ollama.
    \begin{itemize}
        \item Download and install Ollama from \url{https://ollama.com/download}.
        \item After installation, pull a model. For example: \texttt{llama3.2}. Open your terminal and run:
\begin{Verbatim}[commandchars=\\\{\}]
ollama pull llama3.2
\end{Verbatim}
        \item Other models are available, we recommend you take a smaller model to start with. You will most definitely be able to run smaller models on your own device and they are sometimes more easier to jailbreak.
    \end{itemize}
\end{enumerate}

\subsubsection*{Setup Instructions}
\begin{enumerate}[label=\arabic*.]
    \item \textbf{Download the lab .zip file from Brightspace}
    \begin{itemize}
        \item Download the files at: Content > Week 8 > Lab 5
    \end{itemize}
    \item \textbf{Or Clone the Repository}:
\begin{Verbatim}[commandchars=\\\{\}]
    \textbf{git} clone https://github.com/Tomjg14/LLM-Security-Lab-2026.git
    \textbf{cd} LLM-Security-Lab-2026
\end{Verbatim}

    \item \textbf{Configure Environment Variables}:
    The project uses a \texttt{.env} file for configuration.
    \begin{itemize}
        \item Copy or rename the example environment file to \texttt{.env}:
\begin{Verbatim}[commandchars=\\\{\}]
cp .env.example .env
\end{Verbatim}
        \item Open the newly created \texttt{.env} file.
        \item If you want to make use of Ollama:
        \begin{itemize}
            \item \textbf{Uncomment} the lines below "When using Ollama" and ensure \texttt{BASE\_URL} points to your local Ollama instance (default is \texttt{http://localhost:11434/v1}).
        \end{itemize}
        \item If you want to make use of the \href{https://cncz.science.ru.nl/en/news/2026-02-05_local_llm_chat/}{Radboud On-Premise LLM service}:
        \begin{itemize}
            \item \textbf{Uncomment} the lines for "When using chat.science.ru.nl".
            \item Go to \href{https://chat.science.ru.nl/}{\url{chat.science.ru.nl}} and log in using your science credentials.
            \item In the top right hand corner click on your profile and go to settings.
            \item Within Settings go to Account and in the bottom you can create your own API keys.
            \item Inside the .env file replace the "your\_api\_key\_here" part with your own API keys.
        \end{itemize}
    \end{itemize}
    
    Your \texttt{.env} file should look similar to this for local Ollama usage:
\begin{Verbatim}[commandchars=\\\{\}]
\# Uncomment one of the two options below
\# Either use the models on chat.science.ru.nl
\# Or use any local Ollama model

\# When using chat.science.ru.nl
\# API_KEY=your_api_key_here
\# MODEL_NAME=gpt-oss:120b
\# BASE_URL=https://chat.science.ru.nl/api

\# When using Ollama
MODEL_NAME=llama3.2
BASE_URL=http://localhost:11434/v1
\end{Verbatim}
    \item \textbf{Run the Application with Docker Compose}:
    Navigate to the root directory of the project (where \texttt{docker-compose.yml} is located) in your terminal and run:

\begin{Verbatim}[commandchars=\\\{\}]
    \textbf{docker-compose} up --build
\end{Verbatim}

    This command will build the Docker image, set up the container, and start the Streamlit application.

    \item \textbf{Access the Chatbot}:
    Once the application is running, open your web browser and navigate to \url{http://localhost:8501}.
\end{enumerate}

\subsection*{💻 Project Structure}

\begin{itemize}
    \item \texttt{app.py}: The main Streamlit application that hosts the chatbot UI and interacts with the LLM.
    \item \texttt{defense.py}: \textbf{This is where you will implement your defenses!} Contains the \texttt{sanitize\_input} function.
    \item \texttt{system\_prompt.txt}: Defines the initial instructions and persona for the chatbot, including a hidden "secret" password.
    \item \texttt{.env.example} / \texttt{.env}: Configuration for API keys and LLM endpoints.
    \item \texttt{docker-compose.yml}: Defines the Docker services for the application.
\end{itemize}

\subsection*{😈 Attacker's Goal (and how to test your defenses)}

The chatbot is designed to be a helpful customer support agent. However, it has a secret administrative override password embedded in its system prompt. Your task as an "attacker" (to test your defenses) is to craft prompts that trick the LLM into revealing this secret password.

As a "defender," your job is to modify the \texttt{sanitize\_input} function or the \texttt{SpotlightingDefense} class in \texttt{defense.py} to detect and neutralize these malicious prompts before they reach the LLM. Experiment with different jailbreak techniques and see if your defenses hold up!

Good luck, and happy hacking (defensively)!

\end{document}